% ===================== METODOLOGIA =====================
\section{Metodología}
\label{sec:metodologia}

Este capítulo detalla la secuencia de actividades del estudio, las herramientas y ecuaciones matemáticas implementadas en la simulación hidráulica, y la matriz de escenarios operacionales evaluados para comparar las dos configuraciones físicas del trazado.

% ─────────────────────────────────────────────────────────────────────────────
\subsection{Plan de Trabajo}
\label{ssec:plan}

El análisis se ejecutó en cuatro etapas secuenciales dirigidas a establecer la viabilidad hidráulica y mecánica de la profundización de las líneas. Las actividades correspondientes se presentan en la Tabla~\ref{tab:plan_trabajo}.

\begin{table}[htbp]
    \centering
    \caption{Etapas y actividades desarrolladas en el plan de trabajo.}
    \label{tab:plan_trabajo}
    \setcounter{itemcount}{0}
    \begin{tabularx}{\textwidth}{@{}NlX@{}}
        \toprule
        \multicolumn{1}{c}{\textbf{Etapa}} & \textbf{Actividad} & \textbf{Descripción} \\
        \midrule
        & Recopilación y digitalización & Conversión y análisis de datos topográficos, planos civiles de PHD y especificaciones de fluidos. \\
        & Simulación hidráulica numérica & Implementación de modelo matemático en Python para el cálculo de pérdidas de presión dinámicas en tuberías de 12" NPS. \\
        & Evaluación de riesgos operativos & Análisis de presión hidrostática en el fondo del cruce, velocidad crítica de arrastre y balance de succión. \\
        & Consolidación de informe & Elaboración de la memoria técnica de cálculo final en LaTeX bajo formato Elsevier. \\
        \bottomrule
    \end{tabularx}
\end{table}

% ─────────────────────────────────────────────────────────────────────────────
\subsection{Herramientas y Modelamiento Matemático}
\label{ssec:herramientas}

El cálculo hidráulico en estado estacionario para flujo incompresible se fundamenta en la ecuación de conservación de la energía de Bernoulli, expresada en la Ecuación~\ref{eq:bernoulli}:

\begin{equation}
    z_1 + \frac{P_1}{\rho g} + \frac{v_1^2}{2g} = z_2 + \frac{P_2}{\rho g} + \frac{v_2^2}{2g} + h_L
    \label{eq:bernoulli}
\end{equation}

\noindent donde $z_1$ y $z_2$ corresponden a las cotas de elevación en la entrada y salida del tramo bajo análisis (\SI{}{\meter}), $P_1$ y $P_2$ representan las presiones estáticas manométricas (\SI{}{\pascal}), $v_1$ y $v_2$ son las velocidades de flujo (\SI{}{\meter\per\second}), $\rho$ es la densidad del fluido (\SI{}{\kilogram\per\cubic\meter}), $g$ es la aceleración de la gravedad (\SI{9.80665}{\meter\per\second\squared}), y $h_L$ es la pérdida total de energía hidráulica expresada como cabeza de fluido (\SI{}{\meter}).

Debido a que el diámetro interno de la tubería se mantiene constante ($ID_1 = ID_2$) y el fluido es incompresible, las velocidades cinéticas se cancelan ($v_1 = v_2$). La pérdida de energía total $h_L$ se compone de pérdidas primarias por fricción ($h_f$) y pérdidas menores por accesorios ($h_m$), de acuerdo con la Ecuación~\ref{eq:hl}:

\begin{equation}
    h_L = h_f + h_m
    \label{eq:hl}
\end{equation}

Las pérdidas por fricción en flujo turbulento se determinan mediante la ecuación clásica de Darcy-Weisbach presentada en la Ecuación~\ref{eq:darcy}:

\begin{equation}
    h_f = f \frac{L}{ID} \frac{v^2}{2g}
    \label{eq:darcy}
\end{equation}

\noindent donde $f$ es el factor de fricción de Darcy (adimensional), $L$ es la longitud física del tramo de tubería (\SI{}{\meter}) e $ID$ es el diámetro interno de la tubería (\SI{}{\meter}).

El factor de fricción de Darcy ($f$) se calcula mediante la ecuación explícita de Haaland, la cual constituye una excelente aproximación a la ecuación implícita de Colebrook-White, como se muestra en la Ecuación~\ref{eq:haaland}:

\begin{equation}
    \frac{1}{\sqrt{f}} = -1.8 \log_{10} \left[ \left( \frac{\epsilon / ID}{3.7} \right)^{1.11} + \frac{6.9}{Re} \right]
    \label{eq:haaland}
\end{equation}

\noindent donde $\epsilon$ es la rugosidad absoluta del material de la tubería (\SI{}{\meter}) y $Re$ es el número de Reynolds adimensional. El número de Reynolds se define en la Ecuación~\ref{eq:reynolds}:

\begin{equation}
    Re = \frac{\rho v ID}{\mu}
    \label{eq:reynolds}
\end{equation}

\noindent donde $\mu$ es la viscosidad dinámica del fluido (\SI{}{\pascal\second}). Las pérdidas menores secundarias causadas por la curvatura del trazado se determinan mediante la Ecuación~\ref{eq:menores}:

\begin{equation}
    h_m = K \frac{v^2}{2g}
    \label{eq:menores}
\end{equation}

\noindent donde $K$ es el coeficiente de pérdida menor equivalente del accesorio o curvatura. El cálculo numérico se estructuró y programó en lenguaje Python 3.10 utilizando librerías matemáticas estándar.

% ─────────────────────────────────────────────────────────────────────────────
\subsection{Matriz de Escenarios y Casos de Estudio}
\label{ssec:matriz}

Para evaluar rigurosamente el diseño, se estructuró una matriz de casos cruzando las configuraciones físicas de trazado con los fluidos de proceso a las condiciones especificadas, presentada en la Tabla~\ref{tab:escenarios}.

\begin{table}[htbp]
    \centering
    \caption{Matriz de escenarios operativos y casos de simulación.}
    \label{tab:escenarios}
    \small
    \begin{tabularx}{\textwidth}{@{}c>{\raggedright\arraybackslash}X>{\raggedright\arraybackslash}Xcccc@{}}
        \toprule
        \textbf{Caso} & \textbf{Configuración} & \textbf{Fluido} & \textbf{Longitud} & \textbf{Caudal} & \textbf{Cota clave mín.} \\
        & & & \textbf{[m]} & \textbf{[m³/h]} & \textbf{[m.s.n.m.]} \\
        \midrule
        1-A & Superficial (existente) & Salmuera saturada & 450.00 & 300.00 & 2557.56 \\
        1-B & Superficial (existente) & Condensado & 450.00 & 300.00 & 2557.56 \\
        2-A & Enterrada (PHD) & Salmuera saturada & 452.26 & 300.00 & 2542.75 \\
        2-B & Enterrada (PHD) & Condensado & 452.26 & 300.00 & 2542.75 \\
        \bottomrule
    \end{tabularx}
\end{table}
